%% ch03 --- Typing: hints, checkers and where the checker lies
%%
%% Written. The brief that stood here is in tools/chapters.json; what the
%% writing contradicted is recorded in CLAUDE.md under *Resolved questions*.

\chapter{Typing: hints, checkers and where the checker lies}
\label{ch:typing}

\noindent
You have spent a decade with a compiler that refuses. Write \code{int x =
"8080";} and nothing is produced: no assembly, no test run, no deployment.
The type system is not advice, it is a gate, and every habit you have about
types was formed on the far side of it.

Python has the annotations and not the gate. A hint is a claim written in the
source, the interpreter never reads it, and the program runs either way. If
you want something to refuse, you install it, configure it and run it
yourself \dash{} a second compiler you have to invite. This chapter is about
what that changes, and about the places where the checker you invited quietly
agrees with you and is wrong.

It borrows one thing from earlier: an environment that can run a listing,
which is Chapter~\ref{ch:packaging}'s. Nothing else here needs Chapters
\ref{ch:runtime} or~\ref{ch:packaging}, and \pkg{pydantic} appears only as
far as this chapter's argument needs it \dash{} the depth is
Chapters~\ref{ch:services} and~\ref{ch:aitoolkit}.

\section{A hint is a claim}
\label{sec:ch03-claim}
\index{typing!hints as claims}

Read the listing below before you run it, and decide what you expect. Three
calls, each one violating the annotation above it. How many raise?

\pyfile{code/ch03/hint_is_a_claim.py}{Three violated annotations, and what
the interpreter has to say about them}{lst:ch03-claim}

\noindent
None of them. This is what it prints:

\transcript{ch03-hint-is-a-claim}

\noindent
\api{shout} is annotated to take a \code{str} and is handed an \code{int}; it
returns a \code{str} anyway, because an f-string formats whatever it is
given. \api{port\_of} is annotated to return an \code{int} and returns a
\code{str}, and the caller who adds a thousand to it will find out somewhere
else entirely. And the annotations are ordinary data hanging off the function
object, readable and writable at run time, which is how \pkg{pydantic} and
\pkg{FastAPI} read your signatures \dash{} and which also means a hint can be
replaced with a lie after the fact and nothing anywhere notices.

\begin{csbox}
The nearest C\# analogue is not \code{dynamic}. It is a codebase where every
type name is a comment: the names are there, they are accurate most of the
time, and the runtime has never looked at one. Nullable reference types are
the closest thing you have already used \dash{} an annotation the compiler
reasons about and the CLR does not.
\end{csbox}

\begin{trapbox}
\textbf{\code{str} means the value is a string.} It does not. It means
somebody intended a string. A parameter annotated \code{str} can be
\code{None} at run time, can be an \code{int}, can be an open socket, and
nothing in the language will say so. Everything the annotation buys you comes
from a tool you chose to run, and the next section is about inviting one.
\end{trapbox}

\mermaidfig{ch03-two-compilers}{The claim, the run, and the second compiler
that has to be invited}{a hint, the interpreter, and the checker}

\section{The map}
\label{sec:ch03-map}
\index{typing!C\# to Python map}

Most of the vocabulary transfers, and the table below is the part you can
read once and keep. Two rows do not transfer and get the rest of this
section.

\begin{center}
\small
\begin{tabularx}{\linewidth}{@{}lX@{}}
\toprule
\textbf{C\#} & \textbf{Python} \\
\midrule
\code{int}, \code{string}, \code{bool} & \code{int}, \code{str}, \code{bool}
  \\
\code{List<int>}, \code{Dictionary<string, int>} & \code{list[int]},
  \code{dict[str, int]} \\
\code{string?}, nullable reference types & \code{str | None}, a genuine
  union you must narrow \\
\code{IEnumerable<T>}, \code{IReadOnlyList<T>} & \api{Iterable[T]},
  \api{Sequence[T]}, from \pkg{collections.abc} \\
An interface, implemented by declaration & a \api{Protocol}, satisfied by
  shape \\
\code{record}, DTO, POCO & \code{dataclass}, \api{TypedDict},
  \api{NamedTuple}, a \pkg{pydantic} model \\
\code{enum} & \api{Enum}, or \api{Literal} for a closed set of values \\
\code{is} pattern, \code{out var} & \code{isinstance}, or a \api{TypeIs}
  predicate of your own \\
\code{object} & \code{object}, which you must narrow before you use \\
\code{dynamic} & \api{Any}, which is not the same thing \\
\code{where T : IComparable} & a bound on the type parameter \\
\code{in T}, \code{out T} & nothing: variance is inferred \\
\bottomrule
\end{tabularx}
\end{center}

\noindent
Generics are the first row worth stopping at. The modern spelling declares
the parameter on the function or the class, with no \api{TypeVar} and no
import:

\pyregion{code/ch03/the_map.py}{generics}{Two generic classes, and the
variance nobody declared}{lst:ch03-generics}

\noindent
Neither class said \code{out} or \code{in}. The checker reads which positions
the parameter appears in \dash{} \api{Reader} only ever returns a \code{T},
\api{Sink} only ever accepts one \dash{} and infers covariance and
contravariance from that. Both calls in \api{variance} are accepted, and
swapping the two arguments is refused. It is the same idea as the C\#
keywords with the declaration deleted, and the compensation is that you can
no longer get the annotation wrong.

\begin{trapbox}
\textbf{\code{List[int]} and \code{Optional[str]}.} Those are the spellings
from about 2019, they are what a search result will give you, and the
builtin generics replaced them years ago. Write \code{list[int]} and
\code{str | None}. If you are reading a tutorial that imports \api{List}
from \pkg{typing}, assume everything else on the page is the same age.
\end{trapbox}

Narrowing is the second row. \code{str | None} is a union and the checker
will not let you forget it, so a value is genuinely two types until a branch
separates them \dash{} which is stricter than nullable reference types and
much easier to live with. Where a plain \code{isinstance} will not do, you
can write your own predicate, and there are two spellings that differ in one
important way:

\pyregion{code/ch03/the_map.py}{narrowing}{\api{TypeIs} narrows both
branches; \api{TypeGuard} narrows one}{lst:ch03-narrowing}

\noindent
\api{TypeGuard} tells the checker what is true when the predicate returns
true and nothing at all about the other branch, so the second check in
\api{length\_with\_typeguard} is not decoration: delete it and the line stops
type-checking. \api{TypeIs} narrows both branches, which is almost always
what you meant.

\begin{note}
\api{TypeIs} buys that by demanding something \api{TypeGuard} does not: the
type you narrow \emph{to} must be assignable to the type you narrow
\emph{from}. That is why the listing takes a \api{Sequence} rather than a
\code{list}. \code{list} is invariant, so \code{list[str]} is not a
\code{list[object]}, and \code{list[object] -> TypeIs[list[str]]} is
rejected where the \api{Sequence} pair is accepted. The variance three
paragraphs up and the narrowing here are the same rule seen twice.
\end{note}

\begin{exercise}{e03_01_narrow}{Narrow a type, not just answer a question}
Write a predicate that answers whether every item of a sequence is a
\code{str}. The test checks the answers \emph{and} the return annotation,
because a predicate that returns \code{bool} tells the checker nothing: the
run-time answer and the static narrowing are two separate jobs and this
exercise wants both. The starter's signature is the wrong one on purpose.
\end{exercise}

\section{The same DTO, four ways}
\label{sec:ch03-dto}
\index{typing!dataclass, TypedDict, NamedTuple}

A record, a DTO and a POCO are one decision in C\#. In Python they are four,
and the choice is not about taste. Before the listing: the same JSON goes
into all four, and in it the port is the string \code{"8080"}. Write down
which of the four you expect to refuse it.

\pyfile{code/ch03/dto_four_ways.py}{One shape, four ways to declare it}{lst:ch03-dto}

\transcript{ch03-dto-four-ways}

\noindent
One of four. \code{@dataclass} generates \api{\_\_init\_\_}, \api{\_\_eq\_\_}
and \api{\_\_repr\_\_} and validates nothing; \api{TypedDict} is a shape for
a dictionary that already exists, so at run time there is no class and no
instance, only a \code{dict}; \api{NamedTuple} is a real tuple, which is why
it compares equal to an anonymous one. Only the \pkg{pydantic} model looks at
the value, and it is the only one of the four that costs anything to
construct.

Note what the checker did and did not see. The JSON arrives as \api{Any}, so
not one of those four constructions needed a comment to get past
\api{pyright}. Write the same wrong value as a literal, two lines further
down, and it is caught on both lines. Nothing about the error changed; only
where the value came from did.

\begin{trapbox}
\textbf{\code{@dataclass} validates its fields, like a record with a
constructor.} It does not, and this is the single most expensive assumption a
C\# engineer brings to Python. A dataclass is a code generator for the
boilerplate, not a gate. If the values come from outside the process, the
dataclass is the wrong tool on its own and \pkg{pydantic} is the right one
\dash{} which is the whole of \S\ref{sec:ch03-boundary}.
\end{trapbox}

\begin{center}
\small
\begin{tabularx}{\linewidth}{@{}lXl@{}}
\toprule
\textbf{Shape} & \textbf{Reach for it when} & \textbf{Checks?} \\
\midrule
\code{dataclass} & the value is already yours and already right & no \\
\api{TypedDict} & a \code{dict} exists and you want its shape written down &
  no \\
\api{NamedTuple} & you want a tuple with names, cheaply & no \\
\pkg{pydantic} & the value came from outside this process & yes \\
\bottomrule
\end{tabularx}
\end{center}

\section{A Protocol is satisfied by shape}
\label{sec:ch03-protocol}
\index{typing!Protocol and structural typing}

An interface in C\# is a contract signed by the implementer: a class
satisfies \code{IDisposable} because it says so. A \api{Protocol} is a
contract written by the \emph{consumer}, and a class satisfies it by having
the right members.

\pyfile{code/ch03/structural.py}{An interface the implementers never heard
of}{lst:ch03-protocol}

\noindent
\api{Connection} does not import \api{Closable}, and neither does
\api{StringIO}, which shipped with the interpreter years before you wrote
your protocol. Both go through \api{close\_all}. That is the thing structural
typing buys and C\# cannot: you can write the interface for a type you do not
own, after the fact, without an adapter.

The run-time half is weaker than the static half, and knowing exactly how
weak is worth more than the feature. \api{runtime\_checkable} makes
\code{isinstance} work against a protocol, and \code{isinstance} asks only
whether the \emph{attributes are present}. It does not look at signatures. So
\api{Detonator}, whose \api{close} takes a parameter and therefore does not
satisfy the protocol at all, passes the run-time check \dash{} and
\api{pyright} refuses to let that line through quietly, which is why the
listing carries an ignore comment naming the rule. Use \code{isinstance}
against a protocol for a coarse filter; use the checker for the answer.

\begin{exercise}{e03_02_protocol}{An interface the implementers never heard of}
Write a runtime-checkable \api{Closable} protocol and a function that closes
every item that satisfies it, leaves everything else alone and returns how
many it closed. Nothing the test passes in will inherit from your protocol,
including one class from the standard library. That is the exercise.
\end{exercise}

\section{Where the checker lies}
\label{sec:ch03-lies}
\index{typing!Any propagation}

Invite the checker and it will catch a great deal. It will also, in four
common situations, tell you a file is clean when it is not. The listing below
has \val{ch03.lies.functions} functions in it. Every one declares that it
returns an \code{int}. Every one returns a \code{str}.

\pyfile{code/ch03/where_it_lies.py}{Four ways to stop the checker
seeing}{lst:ch03-lies}

\noindent
The first two are \api{Any} arriving without anyone deciding anything.
\api{json.loads} is annotated to return \api{Any}, so everything downstream
of it is \api{Any}, and \api{Any} satisfies any declared type at all
\dash{} that is what the word means in the type system, not a bug in the
tool. A package that ships no type information does the same thing to
everything it hands you. The last two are you telling the checker to stop
looking: \api{cast} generates no code and checks nothing, and
\code{type: ignore} removes an error from the report without removing it from
the program. Both are legitimate and both transfer responsibility to you.

So what do the two checkers this book pins make of that file? This is
generated by \code{code/measure/checkers.py}, which runs both and writes down
what each said:

\transcript{ch03-checkers}

\noindent
\api{pyright}, in strict mode, reports \val{ch03.pyright.errors}.
\api{mypy}, under \code{-{}-strict}, reports \val{ch03.mypy.errors} \dash{}
both of them the two \api{Any} returns, and neither of them the \api{cast} or
the \code{type: ignore}, which are silent by design in both tools. That
difference is not a defect in either. \api{pyright} is applying the type
system, in which \api{Any} is compatible with everything; \api{mypy}'s strict
profile adds an opinion the type system does not have, which is that a
function declaring \code{int} should not quietly return \api{Any}.

\begin{trapbox}
\textbf{\api{Any} is like \code{dynamic}: a thing I can contain.} \api{Any}
propagates. One value of it turns everything computed from it into \api{Any}
too, silently, for as far as it travels, and the report stays clean the whole
way. \code{dynamic} in C\# at least fails loudly at the point of use.
\api{Any} does not fail at all; it stops the checking and the program carries
on being wrong somewhere later.
\end{trapbox}

And the checker is not the only second compiler you can invite. Here is a
function whose annotation is correct, whose checker is happy, and which is
wrong:

\pyfile{code/ch03/defaults.py}{A default argument, evaluated once}{lst:ch03-defaults}

\transcript{ch03-defaults}

\begin{trapbox}
\textbf{\code{def f(items: list[str] = [])} is an optional list parameter.}
The default is evaluated once, when the \code{def} line runs, and that one
list is shared by every call that does not pass its own. The annotation is
right and the function is wrong. Pass \code{None} and build the list inside.
\end{trapbox}

\noindent
Neither pinned checker says anything about that file. The linter does:

\transcript{ch03-three-tools}

\noindent
\api{pyright} finds \val{ch03.defaults.pyright}, \api{mypy} finds
\val{ch03.defaults.mypy} and \pkg{ruff} finds \val{ch03.defaults.ruff}. The
lesson generalises past this one rule: a type checker answers questions about
types, and a good deal of what goes wrong in Python is not a question about
types. Run both.

\begin{projectbox}
This chapter adds no stage to \code{trace-assert} \dash{} the first is
Chapter~\ref{ch:testing} \dash{} but the project has been paying this
chapter's bill since its first commit. Open
\code{code/src/trace\_assert/\_\_init\_\_.py} and read why \api{Event}'s
payload uses a named factory rather than \code{default\_factory=dict}. Under
a strict checker the bare \code{dict} is \code{dict[Unknown, Unknown]}, and
the fix is a function with a return type. The same complaint fired again
while \code{measure/checkers.py} was being written, and was answered with a
single documented \api{cast} \dash{} which is one of the four holes above,
used on purpose and in one place.
\end{projectbox}

\begin{exercise}{e03_03_defaults}{A default that is not what it looks like}
This one starts wrong rather than unfinished: the function is annotated
exactly as you would annotate it in C\#, the checker is happy with it, and
the test fails on the second call. Fix it so that two callers who pass
nothing do not share a list.
\end{exercise}

\section{Which checker, and a configuration to copy}
\label{sec:ch03-config}
\index{typing!pyright and mypy}

The two are not interchangeable and the previous section is why. \api{pyright}
is fast, it is what your editor is probably already running, and its strict
mode is a genuinely high bar. \api{mypy} is slower, is the older and more
widely deployed of the two, and its strict profile carries a handful of
opinions \api{pyright} does not have, of which \code{no-any-return} is the
one this chapter measured.

This book pins both and runs \api{pyright} over everything on every push. The
recommendation that follows from the measurement rather than from taste: run
one of them in CI over the whole tree, and if that one is \api{pyright}, know
that \api{Any} leaking out of a declared return is a class of defect it will
not report. Start here, in \code{pyproject.toml}:

\begin{tomlcode}
[tool.pyright]
typeCheckingMode = "strict"
include = ["src", "tests"]

[tool.mypy]
strict = true
files = ["src"]
\end{tomlcode}

\noindent
Both tools take the interpreter version from the environment they run in, so
pin it beside those blocks if your project needs a different one; this book's
own \code{code/pyproject.toml} is the copyable version of all of it. Two
notes on living with strict mode, both learned on this book's own code.
Turn it on at the start of a project or on one package at a time; turning it
on across an existing codebase produces a number of diagnostics nobody will
read. And when strict mode complains that a type is \emph{partially unknown},
it is almost never being pedantic \dash{} it has found a place where you
believe you have a type and do not.

\section{Validate at the boundary, trust inside}
\label{sec:ch03-boundary}
\index{typing!validate at the boundary}

Everything above collapses into one rule, and it is the only thing in this
chapter worth memorising.

\mermaidfig{ch03-boundary}{One place validates; everywhere after it, the
hints mean what they say}{validate at the boundary}

\noindent
Inside your own code, a hint is worth a great deal: the checker reads it, the
editor completes from it, and a wrong one is caught at the point you write
it. At the edges \dash{} JSON, the environment, a form, a queue message, a
response from a service \dash{} a hint is worth nothing at all, because the
value arrives as \api{Any} and the checker stopped looking one line earlier.

So put the checking where the values come in, once, in a place you can name,
and let everything after it be ordinary typed Python. That is what
\pkg{pydantic} is for and why it is in every dependency list in this field,
and it is why \S\ref{sec:ch03-dto}'s table has a column that answers
\emph{no} three times: the other three shapes are for values that are already
yours. Chapters~\ref{ch:services} and~\ref{ch:aitoolkit} take it seriously;
this chapter only needs you to know where the line is.

\begin{csbox}
You have been following this rule for years without placing it, because
ASP.NET Core places it for you: model binding is the boundary, and a bound
model has been validated before your action method sees it. Nothing in
Python puts the boundary anywhere by default. The rule is the same one; what
is new is that where it goes is now your decision, and a codebase that never
made that decision has boundaries scattered through it.
\end{csbox}

\begin{exercise}{e03_04_boundary}{Validate at the boundary}
Turn a JSON document into a typed object or refuse it. The test hands your
function the kind of JSON a boundary actually delivers \dash{} a port that is
a string, a port that is not a number at all, a document with the key missing
\dash{} and checks that what comes back has the types it claims, or that
nothing comes back. Write it by hand or with \pkg{pydantic}; the test cares
about the behaviour and not the route.
\end{exercise}

\section{What to carry out of this chapter}
\label{sec:ch03-summary}

A hint is a claim the runtime never checks, and a checker is a second
compiler you have to invite, configure and run. Once invited it earns its
place: the map above transfers almost entirely, \api{Protocol} does something
C\# interfaces cannot, and strict mode finds real defects. But it stops
looking at \api{Any}, and the two checkers you might run do not stop in the
same place. Put the validation at the boundary, keep the inside typed, and
treat every \api{cast} and every \code{type: ignore} as a promise you have
now made yourself.
